\documentclass[10pt,a4paper]{amsart}
\usepackage[margin=19mm]{geometry}
\usepackage{amsmath,amssymb,mathtools}
\usepackage[scr=rsfs,cal=euler]{mathalfa}
\usepackage{xfrac}
\usepackage[unicode,hidelinks,bookmarks=false]{hyperref}
\newcommand{\ie}{i.e.\ }
\newcommand{\cf}{cf.\ }
\newcommand{\resp}{respectively\ }
\newcommand{\Subsection}{\subsection}
\providecommand{\nobreakdash}{\nobreak\hbox{-}}
\setlength{\emergencystretch}{2em}
\multlinegap0pt
\usepackage{amssymb}
\usepackage{xcolor}
\newtheorem{lemma}{Lemma}
\newcommand{\R}{\mathbb R}
\title{Counterexample to the second eigenfunction having one zero for a non-local Schr\"{o}dinger operator}
\author{Ben Andrews, Sophie Chen}
\date{Source: arXiv:2507.19016v5 (CC BY 4.0)}
\begin{document}
\maketitle

\noindent\emph{Source and licence.} Counterexample to the second eigenfunction having one zero for a non-local Schr\"{o}dinger operator. Version arXiv:2507.19016v5 (CC BY 4.0), distributed by arXiv under the Creative Commons Attribution 4.0 International licence, http://creativecommons.org/licenses/by/4.0/, as recorded in the arXiv licence metadata for this exact version and shown on the abstract page https://arxiv.org/abs/2507.19016v5. Full licence notice: \texttt{licenses/arxiv-2507.19016-CC-BY-4.0.txt}.\par\smallskip

\noindent\textit{Editorial scope.} Counted unit: the article's lemma lemma\_self\_adj\_res\_frac\_Lapl (source0.tex lines 385--396). The article's complete original proof of it is delivered with it (source0.tex lines 397--569, one \texttt{proof} environment of 173 lines). No mathematics is written, repaired or re-derived: the statement and the proof are the article's own lines, in the article's own notation, and no retained line is substituted or re-flowed.  Retained uncounted as the source-local context the proof needs: lines 294--299; lines 366--372.  Nothing was omitted from any retained span, so the package carries no omission record.  No retained line contains a citation, so the wrapper carries no bibliography and no bibliography entry.  No retained line contains a figure environment, an image-inclusion command or drawing code, so no image is delivered and the package declares no asset dependency.  Editorial formatting only, in addition: an AMS \texttt{amsart} preamble that restates the article's own declarations and macros used by the retained lines, namely the environments \texttt{lemma} on separate counters, together with the standard packages those lines need.  The retained environments print in this document as Lemma 1 in order of appearance, whereas the article prints the same items with the numbering of its own full document.  The complete native source of the article as distributed in the arXiv e-print for this exact version is bundled unchanged as \texttt{sources/182183/source0.tex}.
\begin{lemma}\label{lem_equiv_norms} 
    There exists a constant $c>0$ such that for any function $u$ in $H^s_0(\Omega)$,
    \[
        [u]_{H^s(\R^n)}\geq c\|u\|_{L^2(\Omega)}.
    \] 
\end{lemma}

\begin{align}\label{operator_domain_res_frac_lap}
D((-\Delta)^s_{\mathrm{res}}):=\big\{u\in H^s_0(\Omega):\;&\text{there exists }f\in L^2(\Omega)
\text{ such that} \nonumber\\
&\mathcal{E}_{\mathrm{res}}[u,v]
 =\langle f,v\rangle_{L^2(\Omega)}
 \text{ for every }v\in H^s_0(\Omega)\big\}.
\end{align}

\begin{lemma}\label{lemma_self_adj_res_frac_Lapl}
Let $\Omega$ be an open bounded subset of $\mathbb{R}^n$. The restricted
fractional Laplace operator
\[
    (-\Delta)^s_{\mathrm{res}}
    :
    D((-\Delta)^s_{\mathrm{res}})\subset L^2(\Omega)
    \longrightarrow L^2(\Omega),
\]
with operator domain defined by
\eqref{operator_domain_res_frac_lap}, is positive and self-adjoint.
\end{lemma}

\begin{proof} We first show that $(-\Delta)^s_{\mathrm{res}}$ is bijective. Let
$f\in L^2(\Omega)$ and define
\[
    F_f(v):=\langle f,v\rangle_{L^2(\Omega)},
    \qquad v\in H^s_0(\Omega).
\]
By the Cauchy--Schwarz inequality and
Lemma~\ref{lem_equiv_norms},
\[
    |F_f(v)|
    \leq
    \|f\|_{L^2(\Omega)}
    \|v\|_{L^2(\Omega)}
    \leq
    C\|f\|_{L^2(\Omega)}
    \|v\|_{H^s_0(\Omega)}.
\]
$F_f$ is therefore a bounded linear functional on the Hilbert space
$H^s_0(\Omega)$. By the Riesz representation theorem, there exists a
unique $u\in H^s_0(\Omega)$ such that
\[
    \mathcal E_{\mathrm{res}}[u,v]
    =
    \langle f,v\rangle_{L^2(\Omega)}
    \qquad
    \text{for every }v\in H^s_0(\Omega).
\]
Recalling the definition of $D((-\Delta)^s_{\mathrm{res}})$, this is equivalent to
\[
    u\in D((-\Delta)^s_{\mathrm{res}}),
    \qquad
    (-\Delta)^s_{\mathrm{res}}u=f.
\]
Existence and uniqueness for every $f\in L^2(\Omega)$ show
that $(-\Delta)^s_{\mathrm{res}}$ is bijective.

Next, the operator domain $D((-\Delta)^s_{\mathrm{res}})$ is dense in $L^2(\Omega)$. Indeed,
suppose that $h\in L^2(\Omega)$ is orthogonal to $D((-\Delta)^s_{\mathrm{res}})$.
Since $(-\Delta)^s_{\mathrm{res}}$ is onto, there exists
$u\in D((-\Delta)^s_{\mathrm{res}})$ such that
$(-\Delta)^s_{\mathrm{res}}u=h$. Hence,
\[
    0
    =
    \langle h,u\rangle_{L^2(\Omega)}
    =
    \langle (-\Delta)^s_{\mathrm{res}}u,u\rangle_{L^2(\Omega)}
    =
    \mathcal E_{\mathrm{res}}[u,u]
    =
    \|u\|_{H^s_0(\Omega)}^2.
\]
Consequently $u=0$, so
$h=(-\Delta)^s_{\mathrm{res}}u=0$. It follows that
$D((-\Delta)^s_{\mathrm{res}})$ is dense in $L^2(\Omega)$.

For any $u,v\in D((-\Delta)^s_{\mathrm{res}})$,
\begin{align*}
    \langle (-\Delta)^s_{\mathrm{res}}u,v\rangle_{L^2(\Omega)}
    &=
    \mathcal E_{\mathrm{res}}[u,v] \\
    &=
    \mathcal E_{\mathrm{res}}[v,u] \\
    &=
    \langle u,(-\Delta)^s_{\mathrm{res}}v\rangle_{L^2(\Omega)}.
\end{align*}
Thus $(-\Delta)^s_{\mathrm{res}}$ is symmetric. It is also positive,
since
\[
    \langle (-\Delta)^s_{\mathrm{res}}u,u\rangle_{L^2(\Omega)}
    =
    \mathcal E_{\mathrm{res}}[u,u]
    \geq 0.
\]

It remains to show that $(-\Delta)^s_{\mathrm{res}}$ is self-adjoint.
Since $D((-\Delta)^s_{\mathrm{res}})$ is dense in $L^2(\Omega)$ and symmetric, this implies that  
\[
    D((-\Delta)^s_{\mathrm{res}})
    \subset
    D\bigl({(-\Delta)^s_{\mathrm{res}}}^*\bigr)
\]
and the two operators agree on $D((-\Delta)^s_{\mathrm{res}})$. So we only need to prove the reverse inclusion of the
domains.
% , the adjoint
% ${(-\Delta)^s_{\mathrm{res}}}^*$ is well defined. Indeed, the density
% ensures uniqueness of the vector representing the functional in the
% definition of the adjoint.

% Since $(-\Delta)^s_{\mathrm{res}}$ is symmetric, for every
% $v\in D((-\Delta)^s_{\mathrm{res}})$ we have
% \[
%     \langle (-\Delta)^s_{\mathrm{res}}x,v\rangle_{L^2(\Omega)}
%     =
%     \langle x,(-\Delta)^s_{\mathrm{res}}v\rangle_{L^2(\Omega)}
%     \qquad
%     \text{for every }x\in D((-\Delta)^s_{\mathrm{res}}).
% \]
% Comparing this with the definition of the adjoint shows that
% $v\in D({(-\Delta)^s_{\mathrm{res}}}^*)$ and
% \[
%     {(-\Delta)^s_{\mathrm{res}}}^*v
%     =
%     (-\Delta)^s_{\mathrm{res}}v.
% \]
% Hence
% \[
%     (-\Delta)^s_{\mathrm{res}}
%     \subset
%     {(-\Delta)^s_{\mathrm{res}}}^*,
% \]
% where this operator inclusion means that
Let $v\in D\bigl({(-\Delta)^s_{\mathrm{res}}}^*\bigr)$. As $(-\Delta)^s_{\mathrm{res}}$ is onto $L^2(\Omega)$, there exists
$u\in D((-\Delta)^s_{\mathrm{res}})$ such that
\[
    (-\Delta)^s_{\mathrm{res}}u
    =
    {(-\Delta)^s_{\mathrm{res}}}^*v.
\]
Then, for every $x\in D((-\Delta)^s_{\mathrm{res}})$,
\begin{align*}
    \langle (-\Delta)^s_{\mathrm{res}}x,v\rangle_{L^2(\Omega)}
    &=
    \langle x,{(-\Delta)^s_{\mathrm{res}}}^*v\rangle_{L^2(\Omega)} \\
    &=
    \langle x,(-\Delta)^s_{\mathrm{res}}u\rangle_{L^2(\Omega)} \\
    &=
    \langle (-\Delta)^s_{\mathrm{res}}x,u\rangle_{L^2(\Omega)},
\end{align*}
where the first equality follows from the definition of the adjoint
and the last from symmetry. Hence,
\[
    \langle (-\Delta)^s_{\mathrm{res}}x,v-u\rangle_{L^2(\Omega)}
    =
    0
    \qquad
    \text{for every }x\in D((-\Delta)^s_{\mathrm{res}}).
\]
As $(-\Delta)^s_{\mathrm{res}}$ is onto, the vector
$(-\Delta)^s_{\mathrm{res}}x$ ranges over all of $L^2(\Omega)$.
Therefore, writing $f:=(-\Delta)^s_{\mathrm{res}}x$,
\[
    \langle f,v-u\rangle_{L^2(\Omega)}
    =
    0
    \qquad
    \text{for every }f\in L^2(\Omega).
\]
Taking $f=v-u$ yields
\[
    \|v-u\|_{L^2(\Omega)}^2=0,
\]
and so $v=u\in D((-\Delta)^s_{\mathrm{res}})$. We thus conclude that
\[
    D\bigl({(-\Delta)^s_{\mathrm{res}}}^*\bigr)
    \subset
    D((-\Delta)^s_{\mathrm{res}}).
\]

Together with the reverse inclusion obtained from symmetry,
\[
    D\bigl({(-\Delta)^s_{\mathrm{res}}}^*\bigr)
    =
    D((-\Delta)^s_{\mathrm{res}}),
\]
and the two operators agree on this common domain. Thus,
\[
    (-\Delta)^s_{\mathrm{res}}
    =
    {(-\Delta)^s_{\mathrm{res}}}^*,
\]
showing that $(-\Delta)^s_{\mathrm{res}}$ is self-adjoint.
\end{proof}

\end{document}
